%% kapa: exam=2 round=28 year=2017 period=3 subject=compensation_law subject_ko=감정평가 및 보상법규 seq=4 number=4 points=10
%% source: Q-Net 감정평가사 기출문제 (cst00309) artlSeq=5201712 file=raw/1008/5201712/2017년 감정평가사 2차시험 기출문제(감정평가 및 보상법규).pdf
%% conversion: kordoc 4.24.51 -> tools/gen2.py (auto); review=proofread-sonnet; flags=none
\begin{longquestion}{【문제 4】}
甲 소유의 토지를 포함하는 일단의 토지가 「공공토지의 비축에 관한 법률」에 따라 X읍-Y읍간 도로사업용지 비축사업(이하 ‘이 사건 비축사업’이라 함) 지역으로 지정되었고, 한국토지주택공사를 사업시행자로 하여 2014. 3. 31. 이 사건 비축사업에 대하여 「공익사업을 위한 토지 등의 취득 및 보상에 관한 법률」에 따른 사업인정 고시가 있었다. 한편, 관할 도지사는 X읍-Z읍간 도로확포장공사와 관련하여 2016. 5. 1. 도로구역을 결정·고시하였는데, 甲의 토지는 도로확포장공사가 시행되는 도로구역 인근에 위치하고 있다. 이후 이 사건 비축사업을 위하여 甲 소유의 토지에 대해서 2016. 7. 5. 관할 토지수용위원회의 수용재결이 있었는 바, 위 도로확포장공사로 인하여 상승된 토지가격이 반영되지 않은 감정평가가격으로 보상금이 결정되었다. 이에 甲은 도로확포장공사로 인한 개발이익이 배제된 보상금 결정은 위법하다고 주장하는 바, 甲의 주장이 타당한지에 관하여 설명하시오. (10점) \par
\end{longquestion}
